%!TeX root=../sensetop.tex
\chapter[Chapter \thechapter]{}
\lettrine[lraise=0.3]{E}{linor}'s curiosity to see Mrs Ferrars was satisfied. She had found in her every thing that could tend to make a farther connection between the families undesirable. She had seen enough of her pride, her meanness, and her determined prejudice against herself, to comprehend all the difficulties that must have perplexed the engagement, and retarded the marriage, of Edward and herself, had he been otherwise free; and she had seen almost enough to be thankful for her \textit{own} sake, that one greater obstacle preserved her from suffering under any other of Mrs Ferrars's creation, preserved her from all dependence upon her caprice, or any solicitude for her good opinion. Or at least, if she did not bring herself quite to rejoice in Edward's being fettered to Lucy, she determined, that had Lucy been more amiable, she \textit{ought} to have rejoiced.

She wondered that Lucy's spirits could be so very much elevated by the civility of Mrs Ferrars;—that her interest and her vanity should so very much blind her as to make the attention which seemed only paid her because she was \textit{not Elinor}, appear a compliment to herself—or to allow her to derive encouragement from a preference only given her, because her real situation was unknown. But that it was so, had not only been declared by Lucy's eyes at the time, but was declared over again the next morning more openly, for at her particular desire, Lady Middleton set her down in Berkeley Street on the chance of seeing Elinor alone, to tell her how happy she was.

The chance proved a lucky one, for a message from Mrs Palmer soon after she arrived, carried Mrs Jennings away.

<My dear friend,> cried Lucy, as soon as they were by themselves, <I come to talk to you of my happiness. Could anything be so flattering as Mrs Ferrars's way of treating me yesterday? So exceeding affable as she was! You know how I dreaded the thoughts of seeing her; but the very moment I was introduced, there was such an affability in her behaviour as really should seem to say, she had quite took a fancy to me. Now was not it so? You saw it all; and was not you quite struck with it?>

<She was certainly very civil to you.>

<Civil!—Did you see nothing but only civility?—I saw a vast deal more. Such kindness as fell to the share of nobody but me!—No pride, no hauteur, and your sister just the same—all sweetness and affability!>

Elinor wished to talk of something else, but Lucy still pressed her to own that she had reason for her happiness; and Elinor was obliged to go on.

<Undoubtedly, if they had known your engagement,> said she, <nothing could be more flattering than their treatment of you;—but as that was not the case\longdash>

<I guessed you would say so,>—replied Lucy quickly—<but there was no reason in the world why Mrs Ferrars should seem to like me, if she did not, and her liking me is every thing. You shan't talk me out of my satisfaction. I am sure it will all end well, and there will be no difficulties at all, to what I used to think. Mrs Ferrars is a charming woman, and so is your sister. They are both delightful women, indeed!—I wonder I should never hear you say how agreeable Mrs Dashwood was!>

To this Elinor had no answer to make, and did not attempt any.

<Are you ill, Miss Dashwood?—you seem low—you don't speak;—sure you an't well.>

<I never was in better health.>

<I am glad of it with all my heart; but really you did not look it. I should be sorry to have \textit{you} ill; you, that have been the greatest comfort to me in the world!—Heaven knows what I should have done without your friendship.>

Elinor tried to make a civil answer, though doubting her own success. But it seemed to satisfy Lucy, for she directly replied,

<Indeed I am perfectly convinced of your regard for me, and next to Edward's love, it is the greatest comfort I have. Poor Edward! But now there is one good thing, we shall be able to meet, and meet pretty often, for Lady Middleton's delighted with Mrs Dashwood, so we shall be a good deal in Harley Street, I dare say, and Edward spends half his time with his sister—besides, Lady Middleton and Mrs Ferrars will visit now;—and Mrs Ferrars and your sister were both so good to say more than once, they should always be glad to see me. They are such charming women!—I am sure if ever you tell your sister what I think of her, you cannot speak too high.>

But Elinor would not give her any encouragement to hope that she \textit{should} tell her sister. Lucy continued.

<I am sure I should have seen it in a moment, if Mrs Ferrars had took a dislike to me. If she had only made me a formal courtesy, for instance, without saying a word, and never after had took any notice of me, and never looked at me in a pleasant way—you know what I mean—if I had been treated in that forbidding sort of way, I should have gave it all up in despair. I could not have stood it. For where she \textit{does} dislike, I know it is most violent.>

Elinor was prevented from making any reply to this civil triumph, by the door's being thrown open, the servant's announcing Mr Ferrars, and Edward's immediately walking in.

It was a very awkward moment; and the countenance of each showed that it was so. They all looked exceedingly foolish; and Edward seemed to have as great an inclination to walk out of the room again, as to advance farther into it. The very circumstance, in its unpleasantest form, which they would each have been most anxious to avoid, had fallen on them.—They were not only all three together, but were together without the relief of any other person. The ladies recovered themselves first. It was not Lucy's business to put herself forward, and the appearance of secrecy must still be kept up. She could therefore only \textit{look} her tenderness, and after slightly addressing him, said no more.

But Elinor had more to do; and so anxious was she, for his sake and her own, to do it well, that she forced herself, after a moment's recollection, to welcome him, with a look and manner that were almost easy, and almost open; and another struggle, another effort still improved them. She would not allow the presence of Lucy, nor the consciousness of some injustice towards herself, to deter her from saying that she was happy to see him, and that she had very much regretted being from home, when he called before in Berkeley Street. She would not be frightened from paying him those attentions which, as a friend and almost a relation, were his due, by the observant eyes of Lucy, though she soon perceived them to be narrowly watching her.

Her manners gave some re-assurance to Edward, and he had courage enough to sit down; but his embarrassment still exceeded that of the ladies in a proportion, which the case rendered reasonable, though his sex might make it rare; for his heart had not the indifference of Lucy's, nor could his conscience have quite the ease of Elinor's.

Lucy, with a demure and settled air, seemed determined to make no contribution to the comfort of the others, and would not say a word; and almost every thing that \textit{was} said, proceeded from Elinor, who was obliged to volunteer all the information about her mother's health, their coming to town, \&c. which Edward ought to have inquired about, but never did.

Her exertions did not stop here; for she soon afterwards felt herself so heroically disposed as to determine, under pretence of fetching Marianne, to leave the others by themselves; and she really did it, and \textit{that} in the handsomest manner, for she loitered away several minutes on the landing-place, with the most high-minded fortitude, before she went to her sister. When that was once done, however, it was time for the raptures of Edward to cease; for Marianne's joy hurried her into the drawing-room immediately. Her pleasure in seeing him was like every other of her feelings, strong in itself, and strongly spoken. She met him with a hand that would be taken, and a voice that expressed the affection of a sister.

<Dear Edward!> she cried, <this is a moment of great happiness!—This would almost make amends for every thing!>

Edward tried to return her kindness as it deserved, but before such witnesses he dared not say half what he really felt. Again they all sat down, and for a moment or two all were silent; while Marianne was looking with the most speaking tenderness, sometimes at Edward and sometimes at Elinor, regretting only that their delight in each other should be checked by Lucy's unwelcome presence. Edward was the first to speak, and it was to notice Marianne's altered looks, and express his fear of her not finding London agree with her.

<Oh, don't think of me!> she replied with spirited earnestness, though her eyes were filled with tears as she spoke, <don't think of \textit{my} health. Elinor is well, you see. That must be enough for us both.>

This remark was not calculated to make Edward or Elinor more easy, nor to conciliate the good will of Lucy, who looked up at Marianne with no very benignant expression.

<Do you like London?> said Edward, willing to say any thing that might introduce another subject.

<Not at all. I expected much pleasure in it, but I have found none. The sight of you, Edward, is the only comfort it has afforded; and thank Heaven! you are what you always were!>

She paused—no one spoke.

<I think, Elinor,> she presently added, <we must employ Edward to take care of us in our return to Barton. In a week or two, I suppose, we shall be going; and, I trust, Edward will not be very unwilling to accept the charge.>

Poor Edward muttered something, but what it was, nobody knew, not even himself. But Marianne, who saw his agitation, and could easily trace it to whatever cause best pleased herself, was perfectly satisfied, and soon talked of something else.

<We spent such a day, Edward, in Harley Street yesterday! So dull, so wretchedly dull!—But I have much to say to you on that head, which cannot be said now.>

And with this admirable discretion did she defer the assurance of her finding their mutual relatives more disagreeable than ever, and of her being particularly disgusted with his mother, till they were more in private.

<But why were you not there, Edward?—Why did you not come?>

<I was engaged elsewhere.>

<Engaged! But what was that, when such friends were to be met?>

<Perhaps, Miss Marianne,> cried Lucy, eager to take some revenge on her, <you think young men never stand upon engagements, if they have no mind to keep them, little as well as great.>

Elinor was very angry, but Marianne seemed entirely insensible of the sting; for she calmly replied,

<Not so, indeed; for, seriously speaking, I am very sure that conscience only kept Edward from Harley Street. And I really believe he \textit{has} the most delicate conscience in the world; the most scrupulous in performing every engagement, however minute, and however it may make against his interest or pleasure. He is the most fearful of giving pain, of wounding expectation, and the most incapable of being selfish, of any body I ever saw. Edward, it is so, and I will say it. What! are you never to hear yourself praised!—Then you must be no friend of mine; for those who will accept of my love and esteem, must submit to my open commendation.>

The nature of her commendation, in the present case, however, happened to be particularly ill-suited to the feelings of two thirds of her auditors, and was so very unexhilarating to Edward, that he very soon got up to go away.

<Going so soon!> said Marianne; <my dear Edward, this must not be.>

And drawing him a little aside, she whispered her persuasion that Lucy could not stay much longer. But even this encouragement failed, for he would go; and Lucy, who would have outstaid him, had his visit lasted two hours, soon afterwards went away.

\begin{letter}
\begin{figure}[tbph]
\centering
\includegraphics[width=\linewidth]{35aside}
\caption{Drawing him a little aside}
\end{figure}
\end{letter}
\begin{a4}
\begin{figure}[tbph]
\centering
\includegraphics[width=.9\linewidth]{35aside}
\caption{Drawing him a little aside}
\end{figure}
\end{a4}

<What can bring her here so often?> said Marianne, on her leaving them. <Could not she see that we wanted her gone!—how teazing to Edward!>

<Why so?—we were all his friends, and Lucy has been the longest known to him of any. It is but natural that he should like to see her as well as ourselves.>

Marianne looked at her steadily, and said, <You know, Elinor, that this is a kind of talking which I cannot bear. If you only hope to have your assertion contradicted, as I must suppose to be the case, you ought to recollect that I am the last person in the world to do it. I cannot descend to be tricked out of assurances, that are not really wanted.>

She then left the room; and Elinor dared not follow her to say more, for bound as she was by her promise of secrecy to Lucy, she could give no information that would convince Marianne; and painful as the consequences of her still continuing in an error might be, she was obliged to submit to it. All that she could hope, was that Edward would not often expose her or himself to the distress of hearing Marianne's mistaken warmth, nor to the repetition of any other part of the pain that had attended their recent meeting—and this she had every reason to expect.